%============================================================
\chapter{Nhân Quả -- Gieo Gì Gặt Nấy (Cause and Effect)}
\begin{center}
  \textit{\color{truyenblue}What you send out comes back to you.}\\
  \textit{(Điều bạn gieo đi sẽ quay trở lại với bạn.)}\\[4pt]
  \small\color{truyengold}--- Quy luật muôn đời ---
\end{center}
\ngancach

\section{Tề Hoàn Công Và Cậu Bé Chăn Bò}
\begin{truyen}{Tề Hoàn Công Và Cậu Bé Chăn Bò}{Trung Hoa Cổ Đại}
\chuhoa{T}{}T ề Hoàn Công -- bá chủ thời Xuân Thu -- một lần đi săn lạc đường trong rừng sâu.\\
\textit{(Duke Huan of Qi, overlord of the Spring and Autumn period, once got lost in a deep forest while hunting.)}

Ông gặp một cậu bé chăn bò liền hỏi đường ra kinh thành.\\
\textit{(He met a young cowherd and asked for the way back to the capital.)}

Cậu bé không chỉ chỉ đường mà còn dẫn đường tận nơi, từ chối nhận vàng thưởng.\\
\textit{(The boy not only gave directions but led the way himself, refusing to accept the reward gold.)}

Nhiều năm sau, cậu bé ấy -- Ninh Thích -- vì tài trị quốc được Hoàn Công phong làm đại thần.\\
\textit{(Years later, that boy -- Ning Qi -- was made a minister by Duke Huan for his talent in governance.)}

Ninh Thích cả đời tận trung phục vụ, mang lại thái bình thịnh trị cho cả nước Tề.\\
\textit{(Ning Qi served loyally all his life, bringing peace and prosperity to the entire state of Qi.)}

\end{truyen}
\begin{baihoc}
  Lòng tốt không tính toán của một cậu bé chăn bò đã gieo mầm cho sự nghiệp lớn lao. Gieo tử tế, gặt trọng dụng.\\
  \textit{(The selfless kindness of a young cowherd planted the seed of a great career. Sow kindness, reap recognition.)}
\end{baihoc}

\section{Người Nông Dân Và Con Rắn Đông Cứng}
\begin{truyen}{Người Nông Dân Và Con Rắn Đông Cứng}{Aesop -- Hy Lạp Cổ Đại}
\chuhoa{M}{}M ùa đông băng giá, một người nông dân nhặt được con rắn đông cứng bên đường.\\
\textit{(In the freezing winter, a farmer found a snake frozen stiff by the roadside.)}

Thương hại, ông đặt rắn vào lòng ủ ấm dù mọi người can ngăn.\\
\textit{(Out of pity, he placed the snake in his bosom to warm it despite everyone's warnings.)}

Khi rắn hồi tỉnh và ấm lại, nó lập tức cắn người đã cứu mình.\\
\textit{(When the snake revived and warmed up, it immediately bit the man who had saved it.)}

Trước khi chết, người nông dân thở dài: ``Ta đã gieo lòng tốt nhầm chỗ.''\\
\textit{(Before dying, the farmer sighed: ``I sowed kindness in the wrong place.'')}

\end{truyen}
\begin{baihoc}
  Nhân quả dạy ta không chỉ gieo điều tốt mà còn phải gieo đúng nơi, đúng người, đúng lúc. Trí tuệ phân biệt thiện -- ác là bước đầu của nhân quả sáng suốt.\\
  \textit{(Cause and effect teaches us not only to do good, but to do it in the right place, for the right person, at the right time. The wisdom to distinguish good from evil is the first step of enlightened karma.)}
\end{baihoc}

\section{Rockefeller Trả Ơn Người Lái Tàu}
\begin{truyen}{Rockefeller Trả Ơn Người Lái Tàu}{Nước Mỹ -- Thế Kỷ XIX}
\chuhoa{J}{}J ohn D.\ Rockefeller thuở nhỏ khi còn nghèo một lần bị mắc kẹt giữa bão tuyết.\\
\textit{(When John D. Rockefeller was still young and poor, he was once stranded in a blizzard.)}

Một người lái tàu già không quen biết đã dừng xe chở cậu bé về tận nhà an toàn.\\
\textit{(An old stranger stopped his cart and drove the young boy safely all the way home.)}

Ông lái tàu từ chối mọi lời cảm ơn: ``Tôi chỉ làm điều bình thường thôi.''\\
\textit{(The old driver refused all thanks: ``I was just doing the ordinary thing.'')}

Hai mươi năm sau, Rockefeller đã thành tỷ phú tìm lại người lái tàu đó.\\
\textit{(Twenty years later, the now-billionaire Rockefeller sought out that driver.)}

Ông mua cho cụ một ngôi nhà và trợ cấp đến cuối đời, nói: ``Ông đã gieo hạt mầm tốt bụng; hôm nay tôi chỉ tưới nước.''\\
\textit{(He bought the old man a house and provided for him for life, saying: ``You sowed the seed of kindness; today I merely water it.'')}

\end{truyen}
\begin{baihoc}
  Điều tử tế bạn làm hôm nay -- dù nhỏ bé -- có thể quay lại theo cách kỳ diệu nhất mà bạn không ngờ tới.\\
  \textit{(The kind act you do today -- however small -- may return to you in the most wondrous way you never imagined.)}
\end{baihoc}

\section{Truyền Thuyết Gương Karma Của Người Tây Tạng}
\begin{truyen}{Gương Karma}{Phật Giáo Tây Tạng}
\chuhoa{T}{}T heo giáo lý Phật Giáo Tây Tạng, mỗi con người đang sống trong chiếc gương vô hình.\\
\textit{(According to Tibetan Buddhist teachings, every person lives inside an invisible mirror.)}

Mọi hành động, lời nói, suy nghĩ đều được gương phản chiếu trả lại y chang.\\
\textit{(Every action, word, and thought is reflected back identically by this mirror.)}

Ai gieo giận dữ, nhận lại giận dữ; ai gieo yêu thương, nhận lại yêu thương.\\
\textit{(Those who sow anger receive anger; those who sow love receive love.)}

Điều khác biệt chỉ là thời gian -- đôi khi gương phản chiếu ngay, đôi khi sau nhiều năm.\\
\textit{(The only difference is time -- sometimes the mirror reflects immediately, sometimes after many years.)}

Vị lạt ma dạy: ``Muốn thay đổi điều ngươi nhận được, trước hết hãy thay đổi điều ngươi gieo đi.''\\
\textit{(The lama taught: ``To change what you receive, first change what you sow.'')}

\end{truyen}
\begin{baihoc}
  Cuộc sống phản chiếu lại chính xác những gì ta đặt vào nó. Muốn nhận điều tốt, hãy là người đầu tiên gieo điều tốt.\\
  \textit{(Life reflects back precisely what we put into it. To receive good, be the first to sow good.)}
\end{baihoc}

\section{Bác Sĩ Norman Bethune Ở Trung Quốc}
\begin{truyen}{Norman Bethune Và Phần Thưởng Muộn Màng}{Canada -- Trung Quốc, 1938}
\chuhoa{N}{}N orman Bethune là bác sĩ người Canada bỏ nghề béo bở ở quê nhà sang Trung Quốc giúp đỡ.\\
\textit{(Norman Bethune was a Canadian doctor who left a lucrative career at home to go to China and help.)}

Ông phẫu thuật ngay trên chiến trường không có thuốc gây mê, ánh đèn dầu leo lắt.\\
\textit{(He performed surgery right on the battlefield with no anesthesia, by the light of flickering oil lamps.)}

Nhiều lần bị thương và nhiễm trùng, ông vẫn không rời bỏ thương binh.\\
\textit{(Wounded and infected multiple times, he still refused to abandon the wounded.)}

Ông qua đời năm 1939 vì nhiễm trùng trong khi mổ, nhưng di sản ông để lại vĩ đại hơn nhiều cuộc đời.\\
\textit{(He died in 1939 from an infection during surgery, but the legacy he left was greater than many lifetimes.)}

Mao Trạch Đông viết bài ``Học tập Bạch Cầu Ân'' để cả quốc gia học theo tấm gương hy sinh của ông.\\
\textit{(Mao Zedong wrote ``In Memory of Norman Bethune'' so the entire nation would learn from his example of sacrifice.)}

\end{truyen}
\begin{baihoc}
  Khi ta gieo trọn vẹn sự cống hiến không vụ lợi, hạt giống ấy nảy mầm thành di sản bất tử vượt thời gian.\\
  \textit{(When we sow complete selfless dedication, that seed grows into an immortal legacy that transcends time.)}
\end{baihoc}
